\setcounter{section}{1}
\setcounter{theorem}{5}
\section{\boldmath The class $\PPA_q$}
\label{sec:definitions}

\paragraph{\boldmath Search Problems in $\FNP$ and $\TFNP$.} A search problem
in $\FNP$ is defined by a polynomial time computable relation
$\calR\subseteq \bit^* \times \bit^*$, that is, for every $(x,y)$, it is
possible to decide whether $(x,y) \in \calR$ in $\poly(|x|, |y|)$ time. A
solution to the search problem on input $x$ is a $y$ such that
$|y| = \poly(|x|)$ and $(x,y) \in \calR$. For convenience, define
$\calR(x) \coloneqq \set{y : (x,y) \in \calR}$. A search problem is {\em total}
if for every input $x \in \bit^*$, there exists $y \in \calR(x)$ such that $|y|
\le \poly(|x|)$. $\TFNP$ is the class of all total search problems in $\FNP$.

\paragraph{Reducibility among search problems.} A search problem $\calR_1$ is
{\em Karp-reducible} (or {\em many-one reducible}) to a search problem
$\calR_2$, or $\calR_1 \preceq \calR_2$ for short, if there exist polynomial-time computable functions $f$ and $g$ such that given any instance $x$ of $\calR_1$, $f(x)$ is an instance of $\calR_2$ such that for any $y \in \calR_2(f(x))$, it holds that $g(x, f(x), y) \in \calR_1(x)$.

On the other hand, we say that $\calR_1$ is {\em Turing-reducible} to $\calR_2$, or $\calR_1 \preceq_T \calR_2$ for short, if there exists a polynomial-time oracle Turing machine that on input $x$ to $\calR_1$, makes oracle queries to $\calR_2$, and outputs a $y \in \calR_1(x)$. In this paper, we primarly deal with Karp-reductions, except in Section7.3, where we compare the two different notions of reductions in the context of $\PPA_q$.

\paragraph{\boldmath $\PPA_q$ via complete problems.} We describe several total search problems (parameterized by $q$) that we show to be inter-reducible. $\PPA_q$ is then defined as the set of all search problems reducible to either one of the search problems defined below.

Recall that Boolean circuits take inputs of the form $\bit^n$ and operate using ($\wedge$, $\vee$, $\lnot$) gates. In addition, we'll also consider circuits acting on inputs in $[q]^n$. We interpret the input to be of the form $(\bit^{\ceil{\log q}})^n$, where the circuit will be evaluated only on inputs where each block of $\ceil{\log q}$ bits represents a element in $[q]$. In the case where $q$ is a prime, we could also represent the circuit as $C:\bbF_q^n \to \bbF_q^n$ with arbitrary gates of the form $g: \bbF_q^2 \to \bbF_q$. However, we can simulate any such gate with $\poly(q)$ many $+$ and $\times$ operations (over $\bbF_q$) along with a constant $(1)$ gate. Hence, in the case of prime $q$, we'll assume that such circuits are composed of only $(+, \times, 1)$ gates.\\[1mm]





\noindent
\begin{minipage}{\linewidth}
\begin{definition}
	\label{def:bipartite} \contour{black}{($\bipartite_q$)}\vspace{-5pt}
	\begin{itemize}[leftmargin=2.15cm,itemsep=4pt]
		\item[\slshape Principle:] A bipartite graph with a non-multiple-of-$q$ degree node has another such node.
		\item[\slshape Object:] Bipartite graph $G = (V \cup U, E)$. Designated vertex $v^* \in V$
		\item[\slshape Inputs:] $\triangleright$ $C:\bit^n \to (\bit^n)^k$, with $(\bit^n)^k$ interpreted as a $k$-subset of $\bit^n$\\
	$\triangleright$ $v^* \in \set{0}\times\bit^{n-1}$ (usually $0^n$)
		\item[\slshape Encoding:] $V\coloneqq\set{0}\times\bit^{n-1}$, $U\coloneqq\set{1}\times\bit^{n-1}$,\\
	$E \coloneqq \set{(v,u) : v \in V \cap C(u) \text{ and } u \in U \cap C(v)}$
		\item[\slshape Solutions:] $v^*$ if $\deg(v^*) \equiv 0 \mod{q}$ and\\
	$v \ne v^*$ if $\deg(v) \not\equiv 0 \mod{q}$
	\end{itemize}
\end{definition}
\end{minipage}\vspace{5mm}



\noindent
\begin{minipage}{\linewidth}
\begin{definition}
	\label{def:lonely} \contour{black}{($\lonely_q$)}\vspace{-5pt}
	\begin{itemize}[leftmargin=2.15cm,itemsep=4pt]
		\item[\slshape Principle:] A $q$-dimensional matching on a non-multiple-of-$q$ many vertices has an isolated node.
		\item[\slshape Object:] $q$-dimensional matching $G = (V, E)$. Designated vertices $V^* \subseteq V$ with $|V^*| \le q-1$
		\item[\slshape Inputs:] $\triangleright$ $C:[q]^n \to [q]^n$\\
	$\triangleright$ $V^* \subseteq [q]^n$ with $|V^*| \le q-1$
		\item[\slshape Encoding:] $V\coloneqq[q]^n$. For distinct $v_1, \ldots, v_q$, edge $e\coloneqq\set{v_1, \ldots, v_q}\in E$ if $C(v_i) = v_{i+1}$, $C(v_q) = v_1$
		\item[\slshape Solutions:] $v \in V^*$ if $\deg(v) = 1$ and\\
	$v \notin V^*$ if $\deg(v) = 0$
	\end{itemize}
\end{definition}
\end{minipage}\vspace{5mm}



\noindent
\begin{minipage}{\linewidth}
\begin{definition}
	\label{def:leaf} \contour{black}{($\leaf_q$)}\vspace{-5pt}
	\begin{itemize}[leftmargin=2.15cm,itemsep=4pt]
		\item[\slshape Principle:] A $q$-uniform hypergraph with a non-multiple-of-$q$ degree node has another such node.
		\item[\slshape Object:] $q$-uniform hypergraph $G = (V, E)$. Designated vertex $v^* \in V$
		\item[\slshape Inputs:] $\triangleright$ $C:\bit^n \to (\bit^{nq})^q$, with $(\bit^{nq})^q$ interpreted as $q$ many $q$-subsets of $\bit^n$\\
	$\triangleright$ $v^* \in \bit^n$ (usually $0^n$)
		\item[\slshape Encoding:] $V\coloneqq\bit^n$. For distinct $v_1, \ldots, v_q$, edge $e\coloneqq\set{v_1, \ldots, v_q}\in E$ if $e \in C(v)$ for all $v \in e$
		\item[\slshape Solutions:] $v^*$ if $\deg(v) \equiv 0 \mod{q}$ and\\
	$v \ne v^*$ if $\deg(v) \not\equiv 0 \mod{q}$
	\end{itemize}
\end{definition}
\end{minipage}\vspace{5mm}


\noindent We remark that $\lonely_q$ and $\leaf_q$ are modulo-$q$ analogs of the $\PPA$-complete problems $\lonely$ and $\leaf$ \cite{r33,r5}. We prove the following theorem in \autoref{sec:equivalences}.

\begin{theorem}\label{thm:PPAqEquivalentProblems}
The problems $\bipartiteq$, $\lonelyq$ and $\leafq$ are inter-reducible.
\end{theorem}

\begin{remark}[Simplifications in describing reductions.]\label{rem:simplifications}
	We will use the following simple conventions repeatedly, in order to simplify the descriptions of reductions between different search problems.
	\begin{enumerate}
						\item We will often use ``algorithms'', instead of ``circuits'' to encode our hypergraphs. It is standard to simulate polynomial-time algorithms by polynomial sized circuits.
		\item While our definitions require vertex sets to be of a very special form, e.g. $\bit^n$ or $[q]^n$, it will hugely simplify the description of our reductions to let vertex sets be of arbitrary sizes. This is not a problem as long as the vertex set is efficiently indexable, that is, elements of $V$ must have a $\poly(n)$ length representation and we must have a poly-time computable bijective map $\varphi : V \to [|V|]$, whose inverse is also poly-time computable. We could then use $\varphi$ to interpret the first $|V|$ elements of $\bit^n$ (or $[q]^n$) as vertices in $V$.

		Note that, we need to ensure that no new solutions are introduced in this process. In the case of $\bipartite_q$ or $\leaf_q$, we simply leave the additional vertices isolated and they don't contribute any new solutions. In the case of $\lonely_q$ we need to additionally ensure that $|V| \equiv 0 \mod{q}$, so that we can easily partition the remaining vertices into $q$-uniform hyperedges thereby not introducing any new solutions.
		\item The above simplification gives us that all our problems have an {\em instance-extension property} (cf. \cite{r10}) -- this will be helpful in proving Theorem5.
		\item To simplify our reductions even further, we'll often describe the edges/hyperdges directly instead of specifying how to compute the neighbors of a given vertex. This is only for simplicity and it will be easy to see how to compute the neighbors of any vertex locally.
	\end{enumerate}
\end{remark}

\begin{figure}[t]

	\centering
	\begin{tikzpicture}[scale=0.8, transform shape]
	\def\hgap{2}
	\def\vgap{2.25}
	\draw[fill=Gblue!15, draw=none] (-7,2.4*\vgap) rectangle (1.8,-0.7);
	\draw[fill=Ggreen!15, draw=none] (9,2.4*\vgap) rectangle (1.8,-0.7);
	\draw[fill=orange!15, draw=none] (9,0.3*\vgap) rectangle (1.8,-0.7);
	\node (bip) at (-5.5,2*\vgap) {$\bipartite_p$};
	\node (leaf) at (-5.5,0) {$\leaf_p$};
	\node (leaf-p) at (-5.5,\vgap) {$\leaf'_p$};
	\node (lone) at (-1,0) {$\lonely_p$};
		\node (sbip) at (-1,2*\vgap) {$\sucBipartite_p$};
	\node (twomat) at (-1,\vgap) {$\twoMatchings_p$};

	\def\vgap{1.5}
	\node (combCh) at (5.5,4.5) {$\cChevalley_p$};
	\node (labelCh) at (6,4.5-\vgap) {$\lChevalley_p$};
	\node (Ch) at (6, 4.5-2*\vgap) {$\Chevalley_p$};
	\node (SIS) at (6, 4.5-3*\vgap) {$\sis_p$};

	\path[{Stealth[length=5pt]}-{Stealth[length=5pt]},line width=.4pt,black]
	(bip) edge (leaf-p)
	(leaf-p) edge (leaf)
	(leaf) edge (lone);

	\path[-{Stealth[length=5pt]},line width=.4pt,black]
	(bip) edge (sbip)
	(sbip) edge (twomat)
	(twomat) edge (lone)
	(lone) edge[bend right=0] ([shift={(-1.5,0)}]combCh.south)
	(combCh) edge (sbip)
	(labelCh) edge ([shift={(0.5,0)}]combCh.south)
	(Ch) edge (labelCh)
	(SIS) edge (Ch);
	\end{tikzpicture}
	\caption{Total search problems studied in this work. An arrow
	$A \rightarrow B$ denotes a reduction $A \reducible B$ that we establish.
	Problems in the \textcolor{Gblue}{blue} region are non-natural problems,
	which are all complete for $\PPA_p$. Problems in the
	\textcolor{Ggreen}{green} region are natural problems of which
	$\cChevalley_p$ is the one we show to be $\PPA_p$--complete. The problem
	in the \textcolor{orange}{orange} region is a cryptographically relevant problem.
				}
	\label{fig:ppa-landscape}
\end{figure}
 
